Semaine 1 (Fait)    : Linux, Bash, Scripting, CLI système
Semaine 2 (Fait)    : Git, Branches, Pull Requests, Workflow collaboratif
Semaine 3 (À venir) : Environnements virtuels, Modularité, src-layout, pyproject.toml
─────────────────────────────────────────────────────────────────────────────
Semaine 4 [Jalon 1] : Tests unitaires (pytest), Linters (Ruff) et CI/CD (GitHub Actions)
Semaine 5 [Jalon 2] : Programmation Orientée Objet (POO) industrielle & Exceptions typées
Semaine 6 [Jalon 3a]: I/O, Requêtage HTTP (Requests), Sérialisation & Cache Local
Semaine 7 [Jalon 3b]: Simulation Stochastique Monte-Carlo & Calcul Parallèle (Joblib/MP)
Semaine 8 [Jalon 4a]: Applications Web Décisionnelles & Dataviz (Streamlit, Plotly, Folium)
Semaine 9 [Jalon 4b]: Algorithmique Avancée & Extensions Métier (Patrons de Conception, 2-opt, Numba)
Semaine 10          : Revue de code par les pairs, Profiling/Optimisation, Préparation Soutenance

Séance 4 — Qualité Logicielle, Tests Unitaires (pytest) & Intégration Continue (Jalon 1)Objectif pédagogique : Sécuriser le code par l'automatisation des tests, mesurer la couverture (pytest-cov) et déployer le robot d'intégration continue GitHub Actions.Résonance MANU / Calcul Scientifique : La non-régression et la vérification de convergence numérique (tester si un résidu décroît comme $\mathcal{O}(h^2)$ ou si une propriété d'invariance est respectée) reposent exactement sur cette méthodologie.Découpage du cours-TP (3h00) :La philosophie du test logiciel (30 min) :Pourquoi les assert dispersés dans des scripts sont une impasse.Anatomie d'un test : motif Arrange - Act - Assert (AAA).Types de tests : tests unitaires, tests d'intégration, smoke tests.Le framework pytest en action (45 min) :Découverte automatique des fichiers test_*.py et fonctions test_*().Paramétrisation élégante avec @pytest.mark.parametrize (tester 15 couples entrée/sortie en 4 lignes).Gestion propre des exceptions attendues : pytest.raises(FighterFaintedError).Les fixtures (@pytest.fixture) : instancier des combattants types réutilisables sans duplication de code.Qualité de code statique : Ruff (25 min) :Différence entre analyse statique (linter) et exécution dynamique (tests).Formatage automatique et règles PEP 8 avec ruff check et ruff format.Intégration dans pyproject.toml sous [tool.ruff].Couverture de code (pytest-cov) (20 min) :Notion de couverture de lignes (line coverage) et de branches (branch coverage).Pourquoi $100\%$ ne veut pas dire "zéro bug", mais pourquoi $<70\%$ signifie "code non testé".Échec automatique avec pytest --cov=src --cov-fail-under=70.Pratique dirigée : Création du workflow GitHub Actions (1h00) :Syntaxe YAML : name, on, jobs, steps, runs-on.Déploiement de .github/workflows/ci.yml.Travail en binôme sur GitHub : l'un pousse une régression ou un code mal formaté, observe le passage au rouge $\times$, l'autre soumet la PR corrective et obtient le badge vert $\checkmark$.Séance 5 — Programmation Orientée Objet (POO) Avancée & Exceptions Métier (Jalon 2)Objectif pédagogique : Construire le moteur central du projet (core/models.py, core/engine.py). Remplacer le code spaghetti par une architecture orientée objet robuste, découplée et documentée.Résonance MANU : Modélisation d'opérateurs différentiels abstraits (BaseOperator), maillages (Mesh), conditions aux limites (BoundaryCondition).Découpage du cours-TP (3h00) :Classes Abstraites et Contrats d'Interface (abc) (45 min) :Pourquoi hériter de abc.ABC ? L'interdiction d'instanciation directe.Décorateur @abstractmethod : forcer les classes dérivées à implémenter take_damage(), is_alive(), reset_stats().Le polymorphisme : manipuler une liste de BaseFighter sans connaître leur type concret sous-jacent.Encapsulation et Propriétés Pythoniques (35 min) :Attributs publics, protégés (_hp) et le mythe du privé absolu en Python (__hp et name mangling).Les propriétés @property et setters @hp.setter.Validation des invariants : borner la santé entre $0$ et $\text{PV}_{\max}$, interdire des statistiques négatives.Méthodes Spéciales (Dunder Methods) & Modèles de Données (30 min) :__repr__ vs __str__ : signature technique pour le débogueur vs affichage utilisateur.__eq__ et __hash__ : égalité d'entités basée sur l'identifiant unique.Les structures légères : dataclasses (utilisation pour BattleResult).Hiérarchie d'Exceptions Personnalisées (core/exceptions.py) (30 min) :Pourquoi ne jamais faire except Exception: ou raise ValueError("texte") à l'aveugle.Arborescence d'exceptions : PyArenaException $\rightarrow$ FighterFaintedError, InvalidStatError.Interception ciblée et propagation propre des erreurs.Pratique autonome : Assemblage du moteur core/engine.py (40 min) :Implémentation du duel tour par tour : phase d'initiative, formule de dégâts normalisée, coup critique stochastique, riposte et condition d'arrêt immédiate sur K.-O.Rédaction de la suite de tests unitaires complète dans tests/test_engine.py.Séance 6 — Entrées/Sorties, Requêtage HTTP, Sérialisation & Stratégie de Cache (Jalon 3a)Objectif pédagogique : Concevoir le pipeline de données (io/fetcher.py, io/loader.py). Savoir ingérer des APIs distantes en intégrant la tolérance aux pannes réseau et la mise en cache locale.Résonance MANU : Import de géométries de maillages, lecture de formats scientifiques (HDF5, NetCDF, CSV structurés), découplage chargement / résolution.Découpage du cours-TP (3h00) :Architecture des échanges Web & API REST (35 min) :Protocole HTTP : verbes (GET), codes de statut (200, 404, 500), formats JSON.Utilisation de la bibliothèque requests : requêtes synchrones, passages de paramètres, timeouts.Résilience et Principe du Cache Local (40 min) :Problématique industrielle : limitation de débit (rate-limiting), absence de connexion internet, lenteur réseau.Implémentation du patron Cache-Aside :PlaintextDemande de la créature N°25
   ├── Le fichier data/cache/25.json existe sur disque ?
   │     ├── OUI : Lecture locale immédiate (0 ms, hors-ligne)
   │     └── NON : Requête PokéAPI -> Sérialisation dans data/cache/25.json -> Retour
Sérialisation avec le module standard json (json.dump, json.load).Parsers de Données Locales (io/loader.py) (35 min) :Parsing d'un CSV sans bibliothèque lourde via csv.DictReader pour charger la matrice des types.Lecture de fichiers géographiques GeoJSON (json.load) pour extraire les coordonnées et métadonnées des arènes régionales.Mocking dans les tests unitaires (unittest.mock) (30 min) :Comment tester fetcher.py sans dépendre d'une connexion internet et sans saturer l'API pendant la CI ?Utilisation de unittest.mock.patch pour simuler une réponse HTTP 200 ou une erreur réseau requests.exceptions.ConnectionError.Pratique autonome : Intégration dans PyArena (40 min) :Finalisation de la classe DataFetcher.Validation de la règle du .gitignore pour le cache disque.Tests unitaires avec mocks validés sur la CI GitHub Actions.Séance 7 — Modélisation Stochastique Monte-Carlo & Calcul Parallèle Multi-cœurs (Jalon 3b)Objectif pédagogique : Implémenter le simulateur d'arène probabiliste (simulation/monte_carlo.py). Exploiter les architectures multi-cœurs via la parallélisation de boucles indépendantes (embarrassingly parallel).Résonance MANU : Méthodes de Monte-Carlo pour l'estimation d'incertitude (UQ), parallélisme OpenMP/multiprocessing sur des solveurs de sous-domaines, vectorisation de boucles.Découpage du cours-TP (3h00) :Reproductibilité Stochastique & Générateurs de Nombres Aléatoires (35 min) :La notion de graine (seed) : pourquoi et comment rendre une simulation aléatoire $100\%$ reproductible.La nouvelle API de NumPy : np.random.default_rng(seed) vs l'ancienne méthode globale np.random.seed().Risque de corrélation de graines en environnement distribué.Formalisation Mathématique des Indicateurs de Synthèse (30 min) :Estimateur ponctuel de proportion $\hat{p}$.Intervalle de confiance asymptotique à 95 % (théorème central limite) :$$\mathrm{IC}_{0.95} = \left[ \hat{p} \pm 1.96 \sqrt{\frac{\hat{p}(1-\hat{p})}{N}} \right]$$Moments statistiques de la durée du duel (moyenne, médiane, écart-type).Le verrou global de l'interpréteur (GIL) et les architectures parallèles (40 min) :Pourquoi le multithreading en pur Python (threading) n'accélère pas le calcul numérique (explication limpide du GIL).La solution par multi-processus : créer $P$ interpréteurs Python indépendants ayant chacun leur espace mémoire.Outils standards : multiprocessing.Pool et surcouche ergonomique joblib.Parallel / delayed.Pratique guidée : Parallélisation d'une boucle de simulation (35 min) :Comparaison de temps : boucle séquentielle (for) vs découpage par lots distribués sur les cœurs disponibles (n_jobs=-1).Mesure du temps avec time.perf_counter().Pratique autonome : Implémentation de run_monte_carlo (40 min) :Écriture de la fonction complète dans simulation/monte_carlo.py.Tests unitaires dans tests/test_simulation.py : tester qu'avec une graine donnée, le résultat est déterministe, et que la parallélisation renvoie les mêmes métriques que le mode séquentiel.Séance 8 — Applications Web Décisionnelles & Restitution Visuelle (Jalon 4a)Objectif pédagogique : Construire l'interface graphique Streamlit (app/main.py), intégrer des visualisations interactives (Plotly) et cartographiques (Folium).Résonance MANU : Interfaces web de contrôle et de surveillance pour solveurs EDP, visualisation interactive de champs de pression/température (isocontours Plotly), tableaux de bord de convergence de résidus.Découpage du cours-TP (3h00) :Le Modèle d'Exécution Réactif de Streamlit (40 min) :Comment Streamlit fonctionne : le script se réexécute entièrement de haut en bas à chaque interaction utilisateur.Les composants de base : st.title, st.sidebar, st.selectbox, st.slider, st.columns, st.metric.La gestion de l'état : préserver des calculs lourds ou des choix avec st.session_state.Mise en cache de fonctions lourdes : @st.cache_data pour le chargement des 151/251 combattants.Visualisation Interactive avec Plotly (40 min) :Pourquoi préférer Plotly à Matplotlib en production web : infobulles au survol, zoom interactif, redimensionnement dynamique.Construction des graphiques exigés :Nuage de points interactif (Attaque vs Défense, taille proportionnelle aux PV).Histogramme de distribution des durées de combat.Boîtes à moustaches comparatives par type élémentaire.Cartographie Spatiale avec Folium & Streamlit (35 min) :Principes des couches cartographiques tuilées (tiles).Utilisation de folium.Map, folium.Marker et folium.GeoJson.Intégration dans Streamlit via le composant streamlit-folium.Injection de pop-ups HTML personnalisées affichant le champion local et son sprite.Organisation Multipage (app/pages/) (25 min) :Structuration d'une application professionnelle en plusieurs vues isolées :01_Simulateur_Duel.py02_Catalogue_EDA.py03_Carte_Occitanie.py04_Synergie_Equipe.py (ou extension)Pratique autonome : Connexion UI $\leftrightarrow$ Package PyArena (40 min) :Importer et exécuter run_monte_carlo depuis le dashboard.Afficher en direct les barres de progression et les indicateurs synthétiques.Séance 9 — Algorithmique Avancée & Extensions Métier (Jalon 4b)Objectif pédagogique : Accompagner la mise en œuvre des 10 extensions thématiques. Aborder des patrons de conception logiciels avancés et des algorithmes d'optimisation discrets/numériques.Résonance MANU : Optimisation combinatoire sous contraintes, recherche de racines stochastiques (dichotomie, quasi-Newton), compilation JIT bas niveau (Numba) pour le calcul intensif.Découpage du cours-TP (3h00) :Patrons de Conception (Design Patterns) appliqués au jeu (45 min) :Patron Stratégie (Strategy) : Découpler l'algorithme de décision du combattant (essentiel pour l'Extension 1 Objets, l'Extension 4 Choix d'attaques, l'Extension 5 Statuts).Patron Observateur / Hook : Intercepter un événement de jeu (on_turn_end, on_damage_received).Atelier Méthodologique par Famille d'Extensions (1h00) :Famille Optimisation & Graphes (Ext 2, 3, 9) :Résolution de l'inversion 2-opt et gain Haversine pour le TSP.Algorithme d'énumération de 3-cycles orientés sur matrice d'adjacence tournoi.Famille Performance & Calcul (Ext 7, 8) :Numba : mode @njit(fastmath=True), contrainte de tableaux NumPy 1D contigus, élimination de tout objet Python.Recherche dichotomique discrète convergente sur fonction de winrate monotone stochastique.Famille Chaînes de Markov & Systèmes à États (Ext 1, 5, 6, 10) :Modélisation d'états d'altération, règles d'absorption et modificateurs climatiques globaux.Session de Développement Dirigée en Équipes (1h15) :Chaque groupe travaille sur son fichier src/pyarena/extensions/*.py.Définition et écriture conjointe des tests unitaires stricts dans tests/test_extension.py (vérification des cas limites).Vérification de la conformité du composant Streamlit dédié.Séance 10 — Revue de Code par les Pairs, Audit Final & Préparation des SoutenancesObjectif pédagogique : Préparer le gel des dépôts (main), réaliser une revue de code croisée en conditions réelles, auditer le tableau contractuel CONTRIBUTIONS.md et calibrer la démonstration orale.Découpage du cours-TP (3h00) :L'exercice de la Revue de Code Croisée (Peer Code Review) (1h00) :Le Groupe A examine la dernière PR du Groupe B.Grille de relecture : clarté du nommage, présence des docstrings, couverture des tests unitaires, respect de l'encapsulation.Rédaction de commentaires constructifs sur l'interface GitHub.Audit Technique Pré-Gel (45 min) :Exécution du script enseignant d'évaluation préliminaire devant les étudiants.Vérification que la CI est au vert sur chaque dépôt.Vérification du seuil de couverture ($\ge 70\%$ ou $\ge 85\%$).Vérification de la conformité du fichier CONTRIBUTIONS.md.Cadrage de la Soutenance Orale de 15 minutes (30 min) :Structuration des 8 minutes de présentation : 2 min architecture POO et choix de conception, 3 min démonstration interactive Streamlit, 3 min focus technique sur l'extension.Préparation aux 5-7 minutes de questions individuelles : savoir justifier n'importe quelle ligne de code dont on est l'auteur.Finalisation et Scellement Git (45 min) :Fusion des dernières PRs sur main.Pose du tag final de version (git tag -a v1.0.0 -m "Release finale finale pour évaluation").Gel officiel des dépôts.Pourquoi cette architecture sert parfaitement la mutualisation avec le M1 MANU l'an prochainLe tableau de correspondance ci-dessous montre comment chaque module conçu cette année pour les Sciences des Données transpose directement ses compétences dans le champ de la modélisation et du calcul pour les EDP :Module logiciel PyArenaCompétence InformatiqueApplication directe M1 MANU (Calcul / EDP)src-layout & pyproject.tomlPackaging standardCréation de bibliothèques d'éléments finis réutilisablespytest & GitHub ActionsIntégration continueVérification de l'ordre de convergence numérique ($L^2$, $H^1$)abc.ABC & polymorphismeAbstraction POOOpérateurs différentiels (Laplacien, Convection), solveursExceptions typéesRobustesse logicielleDétection de non-convergence, maillage dégénéré, division par zéroI/O & ParsersIngestion de formatsImport de fichiers de maillages .msh, formats HDF5 / VTKMonte-Carlo & multiprocessingCalcul parallèleQuantification d'incertitudes (UQ), décomposition de domainesNumba / NumPy (Ext 7)Optimisation bas niveauAssemblage vectorisé de matrices de rigidité creusesStreamlit / Folium / PlotlyVisualisationSuivi de résidus de Newton-Raphson, affichage de champs 2D/3D
